\documentclass[11pt, a4paper]{article}

\usepackage[T1]{fontenc}
\usepackage[utf8]{inputenc}
\usepackage{tgpagella}
\usepackage[spanish, es-noshorthands]{babel}
\usepackage{amsmath, amssymb, bm}
\usepackage{tabularx}
\usepackage{booktabs}
\usepackage[most]{tcolorbox}
\usepackage[margin=2.5cm]{geometry}
\usepackage{ragged2e}
\usepackage{fancyhdr}
\usepackage[american]{circuitikz}

\pagestyle{fancy}
\fancyhf{}
\setlength{\headheight}{16pt}
\lhead{\textbf{UCB Sede Tarija} | Ing. Mecatrónica}
\chead{}
\rhead{IMT-121 Circuitos Electrónicos I | Guía Lab RC}
\lfoot{Prof: M.Sc. Ing. Bernardo Quiroga Turdera}
\rfoot{Página \thepage}
\renewcommand{\headrulewidth}{0.4pt}

\definecolor{ucbblue}{RGB}{0, 51, 102}
\definecolor{warnred}{RGB}{200, 0, 0}

\begin{document}

\begin{tcolorbox}[colback=ucbblue!5!white, colframe=ucbblue, title=\textbf{Guía Experimental: Carga y Descarga de un Circuito RC}]
    \centering \textbf{Registro Instrumental y Estimación de Constante de Tiempo}
\end{tcolorbox}

\section*{1. Propósito de la Experiencia}
Contrastar la predicción analítica de un circuito RC de primer orden con la respuesta real observada mediante multímetro y osciloscopio, estimando experimentalmente la constante de tiempo $\tau$.

\section*{2. Consideraciones de Seguridad (¡Leer antes de energizar!)}
\begin{itemize}
    \item \textbf{Polaridad del Capacitor:} Verifique estrictamente la franja negativa del capacitor electrolítico. Conectarlo al revés provocará su destrucción.
    \item \textbf{Tensión Nominal:} Asegúrese de que el capacitor indique \textbf{mínimo 16V}.
    \item \textbf{Descarga Controlada:} \textbf{NUNCA} cortocircuite directamente el capacitor con un cable para descargarlo. Utilice siempre el pulsador hacia la ruta resistiva de descarga ($R_D$) o posicione el conmutador a la referencia de tierra.
\end{itemize}

\section*{3. Circuito a Implementar}
Se utilizará una batería nominal de $9\,\text{V}$ (mida su valor real). El conmutador principal permite alternar la entrada del circuito entre \textbf{Carga} (conectado a batería) y \textbf{Descarga} (conectado a tierra). Adicionalmente, se incluye un pulsador de seguridad para descargar el capacitor rápidamente antes de un nuevo ensayo.

\vspace{0.6cm}
\begin{center}
\begin{circuitikz}[scale=0.85, transform shape, thick]
    % Fuente de voltaje (Batería)
    \draw (0,0) to[battery1, l_={\Large $V_B = 9\,\text{V}$}] (0,4);

    % Terminales del Conmutador
    \draw (0,4) to[short, -o] (2,4) node[above]{\textbf{Carga}};
    \draw (0,0) -- (2,0) to[short, -o] (2,1.5) node[right]{\textbf{Descarga}};

    % Brazo del Switch (SPDT ilustrado)
    \draw[ultra thick, ucbblue] (3.5, 2.75) -- (2,4);
    \draw[dashed, Stealth-Stealth, gray] (2.2, 3.8) arc (70:270:1.2);
    \draw (3.5, 2.75) node[below]{\small Común} to[short, o-] (4.5, 2.75);

    % Resistencia Principal R
    \draw (4.5, 2.75) to[R, l={\Large $R$}] (8, 2.75);

    % Nodo A y Capacitor
    \draw (8, 2.75) node[circle, fill, inner sep=1.5pt, label=above left:{\large $a$}]{};
    \draw (8, 2.75) to[C, l={\Large $C=470\,\mu\text{F}$}, v={$v_C$}] (8,0);

    % Puntas del Osciloscopio (Azul)
    \draw[ucbblue, thick] (8, 2.75) to[short, *-o] (10, 2.75) node[right]{\textbf{Osc. (+) / CH1}};
    \draw[ucbblue, thick] (8, 0) to [short, *-o] (8,-1) to (10, -1) node[right]{\textbf{Osc. (-) / GND}};

    % Malla de Descarga de Seguridad (Rojo, por fuera para no chocar con textos)
    \draw[warnred, thick] (8, 2.75) -- (8, 5.5)
          to[push button, l={Pulsador Reset}] (14, 5.5)
          to[R, l={$R_D$}] (14, 0);

    % Retorno de masa (GND)
    \draw (0,0) -- (14,0);
    \draw (5,0) node[ground]{};

\end{circuitikz}
\end{center}
\vspace{0.6cm}

\begin{tcolorbox}[colback=gray!5, colframe=gray!50, title=\textbf{Dinámica de Rotación (9 Grupos / 5 Osciloscopios)}]
\textit{Mientras su grupo espera el acceso al osciloscopio:}
\begin{enumerate}
    \item Mida y anote los valores reales de sus resistencias y batería en la Hoja de Datos.
    \item Realice todos los cálculos predictivos ($\tau$, voltaje al infinito).
    \item Efectúe la \textbf{medición manual lenta} usando su multímetro y un cronómetro, anotando el voltaje del capacitor cada 5 segundos.
\end{enumerate}
\end{tcolorbox}

\section*{4. Procedimiento en Osciloscopio}
\begin{enumerate}
    \item Asegúrese de que la pinza de masa (GND) del osciloscopio esté en el negativo común de la batería.
    \item Ajuste la base de tiempo considerando que la gráfica completa durará $\approx 5\tau$.
    \item Capture la curva de carga y anote: $t_{inicial}$ (instante de conmutación), voltaje inicial, y voltaje asintótico final.
    \item Estime $\tau_{exp}$ buscando en la pantalla el tiempo en el cual el voltaje ha alcanzado el $\mathbf{63.2\%}$ de la transición total.
\end{enumerate}

\end{document}
